\section{实验与结果}

\subsection{任务、评估方式与证据范围}

本文评价固定实体对上的十选一关系预测。两种评估使用含相同候选选项的题目：直接解码时，模型自行输出关系文本，以 Exact Match（EM）衡量；候选重排时，评估程序逐一代入十个选项，计算 teacher-forced continuation score，再以最高分关系计算 EM、MRR 和 Hits。重排 EM 等于 H@1。选项相同并不意味着答案选择规则相同，两套 EM 不能直接相减来计算训练收益，也不能把重排 EM 当作原始 GRIP 的生成式准确率。

当前具备训练、直接解码和候选重排完整输出、可进入正文主表的 SC-CL 运行位于 NELL23K，使用 Qwen2.5-7B-Instruct 和 seed=2026；测试集包含 4,943 条关系问题。FB15K237-2 的 B1 与历史 InfoNCE-i0 各有 10,000 条候选评分，但 i0 的候选文件、LoRA rank、训练轮数和损失配置不同，不能构成 SC-CL 的跨数据集复现；该结果仅列于附录表~\ref{tab:historical-fb}。WN18RR 已有单种子探索性直接解码测试，但尚无完整候选重排和匹配对照；CoDEx-M 尚无方法训练结果。补充状态见附录~\ref{sec:wn-exploratory}，后续实验范围见附录~\ref{sec:future-protocol}。

\subsection{先核对候选库与训练链路}

NELL23K 的候选库由训练问题构造，写入 3,253 行，每行九个互异关系，\texttt{rows\_failed=0}。冻结 bank 的 wiring 报告确认训练来源、任务问题对应关系和 manifest 记录；SC-CL Stage 2 完成 10 个 epoch、240 个优化步，保存了 adapter、自由生成预测和候选重排输出。这些检查支持\emph{该次运行的来源与执行完整性}，不证明剩余候选均为真实负例。尤其是与持出正例的重叠和开放世界未知事实仍需单独审计。

\subsection{候选排序：有限配对观察与全量结果}

表~\ref{tab:nell-results} 将严格配对的前 200 条问题与只有 SC-CL 的全量闭集结果分开。前 200 条使用相同问题和相同候选列表；B1 的完整测试候选分数尚未完成，故全量一行没有相应的 B1 对照。这样的排列让可比较的差异与单模型绝对水平各自承担不同的证据作用。

\begin{table}[t]
\caption{NELL23K 已观察结果（百分比；单次 seed=2026）。前 200 条为相同问题和候选列表上的 B1/SC-CL 配对比较；full 为仅有 SC-CL 的 4,943 条测试问题。直接解码和候选重排使用相同的十选一题目，但答案选择规则不同。}
\label{tab:nell-results}
\centering\scriptsize
\begin{tabular}{llrrrrr}
\toprule
评估 & 模型 & $n$ & EM/H@1 & MRR & H@3 & H@10 \\
\midrule
闭集 prefix-200 & \bone & 200 & 88.50 & 93.36 & 97.50 & 100.00 \\
闭集 prefix-200 & \scc-9 & 200 & 91.00 & 94.96 & 99.50 & 100.00 \\
闭集 full & \scc-9 & 4,943 & 91.52 & 95.27 & 99.07 & 100.00 \\
\midrule
直接解码 full & \bone & 4,943 & 87.46 & -- & -- & -- \\
直接解码 full & \scc-9 & 4,943 & 87.03 & -- & -- & -- \\
\bottomrule
\end{tabular}
\end{table}

在配对子集上，SC-CL 的 MRR 比 B1 高 1.60 个百分点，候选 EM 高 2.50 个百分点；逐题 top-1 变化为 9 条由错转对、4 条由对转错、173 条共同正确和 14 条共同错误。净变化为五题，同时存在收益与回退。这里使用的是\emph{前} 200 条，而非随机代表样本；单种子结果也没有不确定性估计，因此不能把差异外推到全部测试问题。全量 SC-CL 闭集 MRR 为 95.27\%，描述的是该模型的绝对排序水平；缺少全量 B1 候选评分时，无法计算对应的全量训练增益。

\subsection{直接解码结果与解释边界}

同一批十选一测试问题上，B1 与 SC-CL 的直接解码 EM 分别为 87.46\% 和 87.03\%，描述性差异为 $-0.43$ 个百分点。当前数据没有显示生成能力提升；这一下降是否稳定或是否来自候选训练，尚需同预算、多种子对照检验。SC-CL 的重排 EM 91.52\% 高于自身直接解码 EM 87.03\%，只说明逐候选评分改变了答案选择方式，不能当作训练带来的 4.49 个百分点收益。GRIP 原论文在 NELL23K 上报告 87.74\% 的生成式准确率，但其附录记录 4,952 条测试关系，本地为 4,943 条；该数值只作外部背景，不参与主表的配对差异计算 \citep{feng2025grip}。

\subsection{计算成本与证据缺口}

该次 NELL 运行在训练集上为 644,094 个关系候选生成冻结 B1 分数，用时 19,027 秒，记录的峰值显存为 15.21 GiB。SC-CL Stage 2 处理 12,148,720 个输入 token，用时约 80,870 秒；全测试候选重排用时 3,567 秒。这些是单次运行的绝对成本。由于缺少同硬件、同预算的 Random-9 和 B1 完整训练/评估对照，当前不能计算可靠的单位收益成本比。

全量 B1 闭集评分、同预算 Random-9 和频率匹配随机对照仍缺失，因此目前无法分别估计候选损失与候选构造的作用。持出正例重叠和开放世界未判定也尚未完成完整审计。附录~\ref{sec:future-protocol} 列出后续比较的判定规则；本节的实证范围限于一次可复核运行与前 200 条的描述性配对观察。
